\section{Desafíos abiertos y ruta de investigación}

\subsection{Los desafíos de largo plazo que plantea el curso}
En la parte final del curso, los desafíos se concentran en tres niveles: el nivel cognitivo (razonamiento más confiable), el nivel de sistema (memoria y planificación más robustas) y el nivel organizacional (evaluación y despliegue más confiables).

\begin{figure}[H]
\centering
\includegraphics[width=0.84\textwidth]{slide-06-open-challenges.jpg}
\caption{Página de la presentación en video, aproximadamente en el minuto 01:24:30: las direcciones futuras se concentran en el razonamiento confiable, la gobernanza de la memoria y la planificación de largo alcance}
\end{figure}

\begin{importantbox}{Palancas técnicas clave para los próximos dos años}
\begin{enumerate}
    \item Razonamiento verificable: conectar los resultados del reasoning con verificadores formales;
    \item Gobernanza de la memoria: resolver la propagación de contaminación, los conflictos de versión y la eliminación por caducidad;
    \item Generalización de la planificación: mejorar la capacidad de transferencia entre tareas y reducir el costo de reajuste por tarea.
\end{enumerate}
\end{importantbox}

\begin{knowledgebox}{Por qué es necesario optimizar conjuntamente «Reasoning + Memory + Planning»}
La optimización de un solo punto cae fácilmente en el efecto del eslabón más débil: si el razonamiento mejora pero la memoria se degrada, los errores se amplifican; si la planificación mejora pero el verificador es débil, el sistema se desvía más rápido. El objetivo de la optimización conjunta es una ganancia estable a nivel de sistema, no un máximo local.
\end{knowledgebox}

\begin{warningbox}{Restricciones reales en el despliegue}
Los escenarios empresariales suelen tener presupuestos de latencia estrictos, límites de permisos y requisitos de auditoría. Si un prototipo de investigación ignora estas restricciones, al migrarlo a producción se producirá un retroceso de rendimiento considerable.
\end{warningbox}

\subsection{Resumen de la sección}
La ruta que propone el curso no consiste en perseguir un único SOTA, sino en construir una pila de sistemas de Agent que sea verificable, reversible y escalable. La competitividad de largo plazo proviene de la capacidad de ingeniería de sistemas.

